\scanpage{98}
\mspage{91}
% archive sheet number: 97
Soit
\[
C=\deleted{\unsure{\mathrm{Filt}|A|F}}\underline{M}^{\pi}/\unsure{x|\chi}
\]
la catégorie suivante, les objets sont les triples
\[
(Z,i,\unsure{\alpha|u})\mid Z\in M^{\pi},\quad i\in I,\quad \alpha:Z\longrightarrow X_i,
\]
les flèches
\[
(Z,i,\alpha)\longrightarrow \overset{\unsure{{-}\,\smile|r\,u}}{(Z',i',\alpha')}
\]
sont les couples $(u,\lambda)$ \unsure{avec|où}
\[
u:Z\longrightarrow Z',\qquad \lambda:i\longrightarrow i'
\]
(i.e. on a $i\leq i'$, et $\lambda$ est l'unique flèche de $i$ dans $i'$)
tels que le diagramme
\diagram{\begin{tikzcd}[ampersand replacement=\&]
Z \arrow[r,"\alpha"] \arrow[d,"u"'] \& X_i \arrow[d,"{\text{transition }\unsure{\lambda x_{i',i}|\lambda^{i'}_{,i}|\lambda_{i',i}}}"] \\
Z' \arrow[r,"\alpha'"] \& X_{i'}
\end{tikzcd}}
commute. On peut donc définir $C$ comme le produit fibré dans le carré cartésien
\diagram{\begin{tikzcd}[ampersand replacement=\&]
\& C \arrow[r,"\text{cofibré}"] \arrow[d] \arrow[dl,dashed,"{\unsure{\mathscr F|\mathscr G|\xi}}"'] \& I \arrow[d,"{x:\unsure{X|i}\longmapsto X_i}"] \\
\underline M^{\pi} \& \underline M^{\pi}/E \arrow[l,dashed,"\text{source}"] \arrow[r,"\text{but}","\text{cofibré}"'] \& E
\end{tikzcd}}
\deleted{\unsure{Comme}} La fibre de $C$ en $i\in\deleted{\unsure{J|E}}I$ est $\mathrm{Filt}^{\pi}/X_i$. On a \deleted{\unsure{prouvé}}
% slanted; the curved leader attaches the lower addition at cancelled prouvé
\inserted{Comme b) $\Rightarrow$ \deleted{\unsure{a petite partie|la petite partie}} \unsure{première moitié de a)|première moitié de b)}). que} ($\underline M^{\pi}$ \unsure{est un|est une}
